Una llave ejerce una fuerza $\vec{F} = (0, -20, 0)$ N sobre una tuerca ubicada en el punto $Q = (0.3, 0, 0)$ m respecto al punto de pivote $P$ (en metros). El torque mecánico es el producto cruz $\vec{\tau} = \overrightarrow{PQ} \times \vec{F}$.
\begin{enumerate}
    \item Calcular el vector de torque $\vec{\tau}$.
    \item Determinar la magnitud del torque en N$\cdot$m.
    \item Describir la dirección de rotación que produce el torque usando la regla de la mano derecha.
\end{enumerate}